\documentclass[12pt,a4paper]{article}
\usepackage[utf8]{inputenc}
\usepackage[T1]{fontenc}
\usepackage[spanish,es-noshorthands]{babel}
\usepackage{geometry}
\usepackage{enumitem}
\usepackage{xcolor}
\usepackage{tcolorbox}
\usepackage[hidelinks]{hyperref}

\newcommand{\goldmedal}{1.\textsuperscript{o}}
\newcommand{\silvermedal}{2.\textsuperscript{o}}
\newcommand{\bronzemedal}{3.\textsuperscript{o}}

\geometry{margin=2.5cm}

\title{\textbf{Conversación sobre Web3}\\
\large Profesor CERO \& Kimi AI}
\author{Transcripción de diálogo}
\date{Agosto 2026}

\begin{document}

\maketitle

\tableofcontents
\newpage

%=============================================================================
\section{Pregunta 1: \textquestiondown Web3 es un internet paralelo?}
%=============================================================================

\begin{tcolorbox}[colback=blue!5!white,colframe=blue!75!black,title={\textbf{Pregunta del Profesor CERO}}]
Cre\'{i} que la Web3 era algo as\'{i} como un internet paralelo.
\end{tcolorbox}

\subsection{Respuesta de Kimi}

\textbf{No es un internet paralelo ni una red separada.} Web3 es m\'{a}s bien una \textbf{capa de conceptos y tecnolog\'{i}as} que se construye \textit{sobre el internet que ya existe} (TCP/IP, HTTP, DNS, navegadores... todo eso sigue ah\'{i}).

\begin{table}[h]
\centering
\begin{tabular}{|l|l|l|l|}
\hline
\textbf{} & \textbf{Web1} & \textbf{Web2} & \textbf{Web3} \\
\hline
\textbf{Era} & 1990--2004 & 2004--hoy & Hoy en adelante \\
\hline
\textbf{Lectura/Escritura} & Solo lectura & Lectura + escritura & Lectura + escritura + \textbf{propiedad} \\
\hline
\textbf{Control} & Empresas & Grandes plataformas & Usuarios (te\'{o}ricamente) \\
\hline
\textbf{Ejemplo} & P\'{a}ginas est\'{a}ticas & Redes sociales, YouTube & dApps, DeFi, NFTs, DAOs \\
\hline
\end{tabular}
\caption{Evoluci\'{o}n de la Web}
\end{table}

\subsubsection{Analog\'{i}a f\'{i}sica}
Imagina que el internet es una ciudad:
\begin{itemize}
    \item \textbf{Web1:} Una ciudad de solo bibliotecas. Entras, lees, no interact\'{u}as.
    \item \textbf{Web2:} La ciudad actual: caf\'{e}s, tiendas, cines, pero todo es de empresas grandes. T\'{u} usas los servicios, pero no posees nada.
    \item \textbf{Web3:} La misma ciudad, pero los vecinos pueden ser due\~{n}os de los negocios, votar en su administraci\'{o}n y llevarse sus puntos a otro lugar si no les gusta. \textbf{Es la misma ciudad, pero con reglas diferentes sobre qui\'{e}n posee qu\'{e}.}
\end{itemize}

\subsubsection{La blockchain en todo esto}
La blockchain (Solana, Ethereum, Bitcoin) es la \textbf{infraestructura de confianza} de Web3. Es como el registro de la propiedad de la ciudad, pero distribuido: nadie lo controla por completo, todos pueden verificarlo, y no se puede borrar.

Pero sigues usando:
\begin{itemize}[noitemsep]
    \item Tu navegador (Chrome, Firefox, Brave)
    \item Tu conexión WiFi o 4G
    \item Los mismos cables submarinos y sat\'{e}lites
    \item URLs y dominios web
\end{itemize}

Solo que ahora, adem\'{a}s, interact\'{u}as con \textbf{smart contracts} y \textbf{wallets} (MetaMask, Phantom, etc.) que te conectan a esa capa descentralizada.

\subsubsection{Conexi\'{o}n con proyectos de IoT}
\begin{enumerate}
    \item \textbf{Integridad de datos:} Un sensor ESP32 puede firmar criptogr\'{a}ficamente sus lecturas y enviarlas a una blockchain, garantizando que nadie las manipul\'{o} despu\'{e}s.
    \item \textbf{Micropagos autom\'{a}ticos:} Un nodo IoT vende datos en tiempo real y recibe pagos instant\'{a}neos en cripto por cada lectura.
    \item \textbf{Identidad descentralizada:} Cada dispositivo tiene su propia wallet y su propia identidad en la red, sin depender de un servidor central.
\end{enumerate}


%=============================================================================
\section{Pregunta 2: IoT, firma criptogr\'{a}fica y smart contracts}
%=============================================================================

\begin{tcolorbox}[colback=green!5!white,colframe=green!75!black,title={\textbf{Pregunta del Profesor CERO}}]
Es decir, un dispositivo IoT que recoja informaci\'{o}n puede firmarla criptogr\'{a}ficamente y enviarla a una blockchain. Eso implica generar una wallet para el dispositivo IoT y una identidad que le permitir\'{a} recibir pagos en criptomonedas por cada consulta de sus datos de acuerdo a un smart contract. \textquestiondown Estoy en lo cierto?
\end{tcolorbox}

\subsection{Respuesta de Kimi}

\textbf{Conceptualmente, s\'{i}, est\'{a}s en lo correcto.} Has capturado la esencia del modelo. Pero hay matizaciones t\'{e}cnicas importantes:

\subsubsection{Lo que entendiste bien}
\begin{enumerate}
    \item \textbf{Firma criptogr\'{a}fica:} El ESP32 puede firmar un paquete de datos usando su clave privada. Eso garantiza que los datos provienen de ese dispositivo espec\'{i}fico y que no fueron alterados despu\'{e}s de la medici\'{o}n.
    \item \textbf{Wallet por dispositivo:} Cada nodo IoT puede tener su propio par de claves (p\'{u}blica/privada). La direcci\'{o}n p\'{u}blica es su identidad y su ``cuenta bancaria'' en la red.
    \item \textbf{Smart contract como intermediario programable:} El contrato puede contener la l\'{o}gica de negocio: \textit{``Si alguien quiere estos datos, pague X criptomonedas a la direcci\'{o}n del dispositivo Y antes de desbloquear el acceso''}.
\end{enumerate}

\subsubsection{Matizaciones t\'{e}cnicas importantes}

\paragraph{No se env\'{i}an los datos crudos a la blockchain.}
La blockchain es terriblemente cara para almacenar datos. Lo que suele hacerse es:
\begin{itemize}
    \item El dispositivo env\'{i}a los datos a un \textbf{almacenamiento off-chain} (IPFS, Filecoin, o un servidor tradicional).
    \item A la blockchain solo se env\'{i}a el \textbf{hash criptogr\'{a}fico} de esos datos (una huella digital) junto con la firma del dispositivo.
    \item Eso demuestra que los datos existen y son aut\'{e}nticos, sin pagar una fortuna por almacenamiento.
\end{itemize}

\paragraph{El dispositivo no ``recibe'' el pago en tiempo real como una alerta.}
Las blockchains no notifican a los dispositivos. Lo que ocurre es:
\begin{itemize}
    \item El smart contract actualiza su estado interno: \textit{``La direcci\'{o}n del dispositivo A ahora tiene un saldo de 0.05 ETH''}.
    \item El dispositivo (o su operador) puede consultar ese saldo cuando quiera y ``retirarlo'' (\textit{claim}) a otra wallet.
\end{itemize}

\paragraph{El costo del gas es el cuello de botella.}
Si cada consulta genera una transacci\'{o}n en la blockchain, el costo de la red (gas) puede superar el valor del micropago. Por eso este modelo funciona mejor en:
\begin{itemize}
    \item \textbf{Solana} (transacciones barat\'{i}simas y r\'{a}pidas).
    \item \textbf{Layer 2} (Polygon, Arbitrum, Base).
    \item \textbf{Canales de estado} o \textbf{rollups}, donde se agrupan muchas microtransacciones antes de escribir en la cadena principal.
\end{itemize}

\subsubsection{Arquitectura t\'{i}pica del flujo}

\begin{verbatim}
+-------------+      Firma + Datos       +--------------+
|   ESP32     | -----------------------> |   Gateway/   |
|  (Sensores) |                          |   Oraculo    |
+-------------+                          +------+-------+
                                                |
                    +---------------------------+
                    v
           +-------------+     Hash + Metadata   +-------------+
           |   IPFS      | <--------------------- |  Blockchain |
           |  (Datos)    |                        |  (Smart     |
           +-------------+                        |   Contract) |
                                                 +------+------+
                                                        |
                              +-------------------------+
                              v
                    +-----------------+
                    |  Usuario paga   | --> Smart contract verifica
                    |  por consulta   |     pago y libera acceso
                    +-----------------+     al hash en IPFS
\end{verbatim}

\subsubsection{Conclusi\'{o}n}
Tu intuici\'{o}n es correcta: \textbf{un sensor puede ser un agente econ\'{o}mico aut\'{o}nomo} que vende datos y cobra por ellos sin intermediarios humanos. Esa es exactamente la promesa de Web3 + IoT.


%=============================================================================
\section{Pregunta 3: \textquestiondown Qu\'{e} tipo de datos se pueden vender en Web3?}
%=============================================================================

\begin{tcolorbox}[colback=orange!5!white,colframe=orange!75!black,title={\textbf{Pregunta del Profesor CERO}}]
\textquestiondown Qu\'{e} tipo de datos, seg\'{u}n tu, se pueden vender en la Web3?
\end{tcolorbox}

\subsection{Respuesta de Kimi}

Se distinguen dos planos: \textbf{lo que ya se vende hoy} (mercados reales) y \textbf{lo que es te\'{o}ricamente posible} pero a\'{u}n no tiene liquidez.

\subsubsection{1. Datos de sensores e IoT}

\textbf{Ejemplos concretos:}
\begin{itemize}
    \item \textbf{Datos ambientales:} Temperatura, humedad, calidad del aire, radiaci\'{o}n UV, niveles de CO$_{2}$. Empresas de seguros, agricultura de precisi\'{o}n y gobiernos locales pagan por esto.
    \item \textbf{Datos de tr\'{a}fico/movilidad:} Flujo vehicular, ocupaci\'{o}n de estacionamientos, condiciones de carreteras.
    \item \textbf{Datos s\'{i}smicos:} Redes de aceler\'{o}metros distribuidos. Invaluables para modelos de alerta temprana.
    \item \textbf{Datos de salud remotos:} Ritmo card\'{i}aco, oximetr\'{i}a, glucosa (siempre anonimizados y con consentimiento).
\end{itemize}

\textbf{Mercado real:} Proyectos como \textbf{Helium} (red LoRaWAN descentralizada) y \textbf{DIMO} (datos de veh\'{i}culos) ya monetizan esto.

\textbf{Ventaja del perfil del profesor CERO:} Quien domine el nodo f\'{i}sico tiene una barrera de entrada real. La mayor\'{i}a de los proyectos Web3 se quedan en software.

\subsubsection{2. Datos para entrenar inteligencia artificial}

\textbf{El problema:} Las grandes empresas de IA entrenan sus modelos con internet sin compensar a los creadores de datos.

\textbf{La propuesta Web3:}
\begin{itemize}
    \item Textos, im\'{a}genes, c\'{o}digo: Artistas, escritores y programadores venden sus datasets directamente.
    \item Datos etiquetados: Empresas de IA pagan por datasets curados y verificados.
    \item Datos de comportamiento humano real: Interacciones, preferencias, decisiones de compra (con consentimiento expl\'{i}cito).
\end{itemize}

\textbf{Mercado real:} Plataformas como \textbf{Ocean Protocol} y \textbf{SingularityNET} intentan crear mercados de datos para IA.

\subsubsection{3. Datos financieros y de mercado}

\textbf{Ejemplos:}
\begin{itemize}
    \item \textbf{Feeds de precios en tiempo real:} Criptomonedas, acciones, commodities. Los or\'{a}culos descentralizados (Chainlink) venden estos datos a smart contracts.
    \item \textbf{Datos on-chain:} An\'{a}lisis de flujos de wallets, comportamiento de traders.
    \item \textbf{Ratings de riesgo crediticio descentralizado:} Tu historial de transacciones en DeFi como ``score crediticio'' para pr\'{e}stamos sin bancos.
\end{itemize}

\textbf{Mercado real:} Este es el m\'{a}s maduro. Los or\'{a}culos son infraestructura cr\'{i}tica de DeFi.

\subsubsection{4. Identidad y reputaci\'{o}n verificable}

\textbf{La idea:} En lugar de que Google o Facebook posean tu identidad digital, t\'{u} la controlas y la vendes (o alquilas) por pedacitos.

\textbf{Ejemplos:}
\begin{itemize}
    \item \textbf{Prueba de humanidad:} Demostrar que eres una persona real (no un bot) sin revelar tu nombre.
    \item \textbf{Credenciales verificables:} Diplomas, certificaciones, historial laboral.
    \item \textbf{Reputaci\'{o}n cross-plataforma:} Tu reputaci\'{o}n como vendedor en eBay, tu karma en Reddit, tus contribuciones en GitHub.
\end{itemize}

\textbf{Mercado real:} A\'{u}n muy te\'{o}rico. La interoperabilidad entre plataformas es un problema t\'{e}cnico y pol\'{i}tico enorme.

\subsubsection{5. Contenido creativo y cultural}

\textbf{Ejemplos:}
\begin{itemize}
    \item M\'{u}sica con derechos de uso tokenizados: stems, metadatos de producci\'{o}n, datos de audiencia.
    \item Datos de storytelling: Universos narrativos, personajes, lore.
    \item \textbf{Datos de dise\~{n}o 3D:} Modelos CAD param\'{e}tricos que otros pueden modificar y usar bajo licencia smart contract.
\end{itemize}

\subsubsection{Lo que NO deber\'{i}a venderse (o al menos no sin regulaci\'{o}n)}

\begin{itemize}
    \item \textbf{Datos biom\'{e}tricos sin consentimiento:} Huellas dactilares, iris, ADN. Una vez en la blockchain, son inmutables e inborrables.
    \item \textbf{Datos de vigilancia masiva:} C\'{a}maras de seguridad, reconocimiento facial en espacios p\'{u}blicos.
    \item \textbf{Datos m\'{e}dicos sensibles:} Aunque anonimizados, la reidentificaci\'{o}n es t\'{e}cnicamente posible con datasets suficientemente grandes.
\end{itemize}

\subsubsection{Opini\'{o}n honesta: ranking de probabilidad de \'{e}xito}

\begin{table}[h]
\centering
\begin{tabular}{|c|l|l|}
\hline
\textbf{Ranking} & \textbf{Tipo de dato} & \textbf{Por qu\'{e}} \\
\hline
\textbf{\goldmedal} & Datos de sensores IoT & Escasez real, verificabilidad t\'{e}cnica, demanda de industrias \\
\hline
\textbf{\silvermedal} & Datos de entrenamiento para IA & Demanda masiva, los grandes modelos necesitan datos frescos \\
\hline
\textbf{\bronzemedal} & Feeds financieros/or\'{a}culos & Infraestructura cr\'{i}tica, ya funciona \\
\hline
\end{tabular}
\caption{Datos con mayor probabilidad de \'{e}xito comercial en Web3}
\end{table}

Los datos personales y de identidad son \textbf{m\'{a}s valiosos \'{e}ticamente que econ\'{o}micamente} por ahora. El mercado a\'{u}n no ha resuelto c\'{o}mo valorar la privacidad.


%=============================================================================
\section*{Cierre}
%=============================================================================

\textit{``Un sensor puede ser un agente econ\'{o}mico aut\'{o}nomo que vende datos y cobra por ellos sin intermediarios humanos. Esa es exactamente la promesa de Web3 + IoT.''}

\vspace{1cm}
\begin{flushright}
\textbf{Kimi AI} --- Conversaci\'{o}n con el Profesor CERO\\
Agosto 2026
\end{flushright}

\end{document}comp